\subsection*{增长速度为何取决于已有数量}
设一个量每时每刻都按当前大小的一定比例增长。已有量越多，同样短时间里增加越多，数学关系是 $N'=rN$。这与每次增加固定数量不同；后者形成直线，前者的增加量自身也在增长。

对于 $a^x$，差商可拆成 $a^x(a^h-1)/h$。当前位置的因子 $a^x$ 被提了出来，其余极限只取决于底数。不同底数给不同的比例常数；选择使该常数等于一的底数，就得到自然底数 $\mathrm e$。它特殊在于局部变化规则最简洁。

\subsection*{连续复利是一条通向同一对象的路}
把单位时间分成 $n$ 小段，每段按 $1/n$ 的比例增长，得到 $(1+1/n)^n$。分段越来越细时，它趋向 $\mathrm e$。有限分段仍有自身的数值，不能把每次“增长百分之一”的离散模型直接当成连续模型。

要把一般正底数纳入同一语言，写 $a^x=\mathrm e^{x\ln a}$。内层把时间刻度缩放 $\ln a$ 倍，链式法则于是给出导数 $a^x\ln a$。额外因子可以理解为增长时钟的快慢，而不是一项无来由的常数。

例如希望一个量一单位时间翻倍，连续增长率应取 $\ln2$。若取 $r=1$，一单位时间的倍率是 $\mathrm e$，并非二。练习时同时检查数量单位与时间单位；将来面对微分方程，指数函数将作为“变化率与自身成比例”的基本解再次出现。
